% Introduction — constraints on H0 from EMRI dark sirens with LISA using the MBH mass

The Hubble constant $\HH$ sets the present-day expansion rate of the
Universe and anchors the absolute distance scale in cosmology.  Two
mature, independent determinations of $\HH$ are now in significant
tension.  Observations of the cosmic microwave background (CMB) by the
\textit{Planck} satellite, interpreted within a flat $\Lambda$CDM
cosmology, yield $\HH = 67.4 \pm 0.5\;\mathrm{km\,s^{-1}\,Mpc^{-1}}$~\cite{Planck:2018vyg}.
The local distance ladder, calibrated through Cepheid variables and
Type~Ia supernovae by the SH0ES collaboration, gives $\HH = 73.0 \pm
1.0\;\mathrm{km\,s^{-1}\,Mpc^{-1}}$~\cite{Riess:2021jrx}, a
discrepancy exceeding $5\sigma$.  Whether this tension reflects
unaccounted systematics in one or both measurements, or signals new
physics beyond $\Lambda$CDM, remains an open question.  Resolving it
demands independent probes of $\HH$ that do not share the systematic
uncertainties of either the CMB or the distance ladder.

Gravitational-wave (GW) observations provide such an independent
probe.  As recognized by Schutz~\cite{Schutz:1986gp}, a coalescing
compact binary is a \textit{standard siren}: the GW signal encodes the
luminosity distance $\dL$ directly, without any need for a cosmic
distance ladder.  If the redshift $z$ of the source is also known, the
pair $(\dL, z)$ constrains $\HH$ through the distance--redshift
relation.  The first demonstration came with GW170817, the binary
neutron star merger detected by LIGO and Virgo, whose electromagnetic
counterpart identified the host galaxy NGC~4993 and yielded $\HH =
70.0^{+12.0}_{-8.0}\;\mathrm{km\,s^{-1}\,Mpc^{-1}}$~\cite{Abbott:2017xzu}.
With a growing population of bright sirens accompanied by
electromagnetic counterparts, percent-level $\HH$ measurements become
feasible within a decade of advanced detector
operations~\cite{Chen:2017rfc}.

Most compact binary mergers, however, lack an identifiable
electromagnetic counterpart.  These \textit{dark sirens} require a
statistical approach to obtain redshift information.  In the galaxy
catalog method~\cite{Schutz:1986gp,Gray:2019ksv}, one cross-matches
the GW sky localization and distance estimate with a galaxy catalog,
identifying candidate host galaxies within the localization volume.
The likelihood for $\HH$ is then constructed by marginalizing over all
potential hosts, weighted by their probability of being the true
source.  As the number of dark siren events grows, the contributions
from incorrect host galaxies average out, and the posterior on $\HH$
converges to the true value.  A key subtlety is that galaxy catalogs
are incomplete beyond a redshift that depends on the survey depth; a
completeness correction must be included to avoid biasing the $\HH$
inference~\cite{Gray:2019ksv,Gair:2022zsa}.

The Laser Interferometer Space Antenna (LISA)~\cite{LISA:2024hlh}, an
ESA-led space-based gravitational-wave observatory targeting the
millihertz band, is expected to launch in the mid-2030s.  Among its
richest source classes are extreme mass-ratio inspirals (EMRIs): a
stellar-mass compact object ($\mu \sim 10\,M_\odot$) spiraling into a
massive black hole ($\Mbh \sim 10^5$--$10^7\,M_\odot$) in the center
of a galaxy~\cite{Babak:2017tow,LISA:2024hlh}.  An EMRI signal
persists in the LISA band for $\sim\!10^5$ orbital cycles over
months to years, enabling extraordinarily precise parameter estimation.
Sky localization uncertainties of $\mathcal{O}(1\;\mathrm{deg}^2)$
and fractional luminosity distance errors of a few percent are
achievable for high signal-to-noise ratio (SNR) events~\cite{Babak:2017tow}.
Because EMRIs are not expected to produce detectable electromagnetic
radiation, they are natural dark siren sources.  Their superior sky
localization compared to ground-based dark sirens, however, dramatically
reduces the number of candidate host galaxies per event, making the
statistical galaxy catalog method particularly powerful.

Laghi et al.~\cite{Laghi:2021pqk} performed the first comprehensive
study of $\HH$ inference from EMRI dark sirens with LISA, using the
galaxy catalog approach.  Their analysis adopted a simplified treatment
of EMRI parameter estimation errors, assigning a fixed fractional
luminosity distance uncertainty of $\sigma(\dL)/\dL = 10\%$ to all
detected events, and used only the GW-measured sky position and
distance to cross-match with galaxy catalogs.  Under these assumptions,
they projected $\HH$ constraints at the $\sim\!4\%$ level from a
four-year LISA mission.

In this work, we extend the EMRI dark siren framework by exploiting an
additional observable that LISA measures with remarkable precision: the
redshifted mass of the central massive black hole, $\Mz = \Mbh(1+z)$.
LISA determines $\Mz$ from the EMRI waveform phase evolution to a
fractional accuracy of $\sigma(\Mz)/\Mz \sim 10^{-5}$--$10^{-4}$,
far exceeding the precision of any electromagnetic mass
proxy~\cite{Babak:2017tow}.  Galaxy catalogs, meanwhile, provide
stellar masses $\Mstellar$ for their entries, which can be mapped to
black hole masses through empirical scaling
relations~\cite{Reines:2015moa}.  The comparison between the
GW-measured $\Mz/(1+z)$ and the catalog-predicted $\Mbh$ for each
candidate galaxy provides a powerful additional discriminant that was
not exploited in Ref.~\cite{Laghi:2021pqk}.

We present a complete EMRI dark siren inference pipeline that
incorporates Fisher-matrix-based Cram\'er--Rao bounds for all
14~EMRI parameters, the GLADE+ galaxy catalog~\cite{Dalya:2021ewn},
a completeness correction following Gray et
al.~\cite{Gray:2019ksv}, and a detection probability model
accounting for EMRI population and selection effects.
We perform the first direct comparison of two likelihood
configurations: a baseline ``without MBH mass'' channel that uses only
the GW-measured luminosity distance and sky position $(\dL, \varphi,
\theta)$ to identify candidate hosts, and an extended ``with MBH
mass'' channel that additionally incorporates the redshifted mass
$\Mz$.  We find that including the MBH mass information tightens the
$\HH$ constraint from $\sim\!4\%$ to $\sim\!1.8\%$ fractional
precision, demonstrating that the mass channel substantially enhances
the dark siren method for EMRIs.

The remainder of this paper is organized as follows.
Section~\ref{sec:method} describes the signal model, parameter
estimation framework, dark siren likelihood construction, galaxy
catalog treatment, and detection probability model.
Section~\ref{sec:results} presents our main results, including
detection statistics, the $\HH$ posterior comparison between the two
likelihood channels, and an assessment of systematic effects.
Section~\ref{sec:discussion} discusses the implications and
limitations of our findings.  We conclude in
Sec.~\ref{sec:conclusions}.  Appendix~\ref{app:parameters} tabulates
the EMRI parameter space used in the simulations.
Throughout, we adopt geometrized units with $G = c = 1$ unless
otherwise noted, and quote $\HH$ in units of
$\mathrm{km\,s^{-1}\,Mpc^{-1}}$.
